% ============================================================
% Capítulo: Trabajos Relacionados (Bloque B.6)
% ============================================================
\chapter{Trabajos relacionados}
\label{cap:trabajos-relacionados}

Este capítulo sitúa a SGB-SaaS frente a la literatura académica indexada
disponible sobre sistemas de gestión bibliotecaria y sobre los patrones
arquitectónicos y de proceso que el sistema efectivamente utiliza. Sigue
una metodología \textbf{adaptada} de PRISMA~2020 -- adaptada porque, como
se explica en la \autoref{sec:estrategia-busqueda}, no es una revisión
sistemática completa ejecutada desde cero para este capítulo, sino un
acervo bibliográfico construido en dos momentos del proyecto y
consolidado aquí; se declara esto explícitamente en vez de presentar un
embudo de selección más grande o más formal del que realmente se
ejecutó.

\section{Estrategia de búsqueda}
\label{sec:estrategia-busqueda}

\textbf{Origen del acervo.} Las referencias de este capítulo provienen de
dos fuentes: (1)~un conjunto de 27 referencias curadas por el equipo en
fases anteriores del proyecto a partir de búsquedas en IEEE~Xplore,
ACM~Digital Library, Scopus, ScienceDirect (Elsevier) y SpringerLink; y
(2)~2 referencias adicionales (\citet{colley2018ormperformance} y
\citet{kusuma2017rediscaching}) identificadas mediante una búsqueda
dirigida realizada específicamente para este capítulo -- vía Crossref e
IEEE~Xplore -- para cubrir dos vacíos temáticos puntuales que el acervo
original no cubría: rendimiento comparado de arquitecturas ORM frente a
SQL optimizado, y el patrón \emph{cache-aside} con Redis. Las 30
referencias del archivo \texttt{docs/bibliografia.bib} se verificaron
una por una contra el registro oficial de Crossref de su DOI antes de
transcribirse; esa verificación encontró 5 discrepancias entre los datos
originalmente reunidos por el equipo y el registro oficial (autores
atribuidos incorrectamente en 3 casos, un DOI que en realidad apunta a
un artículo distinto del atribuido en un caso, y un año/autoría
equivocados en el caso de la referencia NIST). Cada discrepancia se
documenta con su evidencia en el encabezado de \texttt{bibliografia.bib}
y se usó, en todos los casos, el dato verificado por Crossref.

\textbf{Cadena de búsqueda de referencia} (la usada en las búsquedas
originales del equipo y replicada en la búsqueda dirigida de este
capítulo):

\begin{quote}
\ttfamily
("library management system" OR "biblioteca" OR "library automation")\\
AND ("QR code" OR RFID OR chatbot OR "recommendation system" OR\\
\phantom{AND (}"generative AI" OR "stored procedure" OR "cache-aside" OR Redis OR ORM)\\
AND (performance OR empirical OR evaluation OR "case study")
\end{quote}

\textbf{Ventana temporal.} 2017--2026, coherente con el rango real de
publicación de las referencias reunidas (\citet{kusuma2017rediscaching},
2017, es la más antigua; \citet{ayinde2026ai}, 2026, la más reciente).

\textbf{Criterios de inclusión.} (a)~publicación arbitrada por pares en
una revista indexada (Scopus/JCR) o en actas de una conferencia
publicada por IEEE/ACM; (b)~publicada entre 2017 y 2026; (c)~aborda
directamente un sistema de gestión bibliotecaria, o un patrón
arquitectónico/de proceso que SGB-SaaS implementa (ORM/SP, caché Redis,
IA generativa, elicitación de requisitos, prácticas ágiles, diseño
responsivo, TLS); (d)~disponible en inglés con resumen consultable.

\textbf{Criterios de exclusión.} Blogs, tutoriales de código sin
evaluación empírica, contenido de marketing de proveedores, y fuentes
autopublicadas sin arbitraje por pares. Una excepción declarada
explícitamente: \citet{brown2023c4model} (el modelo~C4) se cita en este
documento como referencia normativa de notación arquitectónica -- no
como evidencia de investigación empírica -- y por eso queda fuera del
acervo sometido a cribado de este capítulo.

\section{Proceso de selección}
\label{sec:proceso-seleccion}

El acervo de 30 referencias se dividió en dos grupos con tratamiento
distinto, porque no todas compiten con SGB-SaaS como sistemas
alternativos comparables: solo el primer grupo se sometió al cribado
tipo PRISMA de la \autoref{fig:prisma-flow}.

\textbf{Grupo A -- dominio bibliotecario y arquitectura de acceso a
datos comparable (15 referencias).} Sistemas o estudios de rendimiento
directamente comparables con SGB-SaaS como aplicación:
\citet{ayinde2026ai, liabor2023cataloging, shivaji2024intelligentlms,
rajacic2025gdpr, supriyono2018qr, ng2024libbot, mukund2021rfid,
adetayo2023chatgpt, yan2023chatbotsslr, ehrenpreis2022chatbot,
wayesa2023recsys, hu2025bigdata, ayemowa2024genai,
colley2018ormperformance, kusuma2017rediscaching}.

\textbf{Grupo B -- literatura de apoyo metodológico y arquitectónico (15
referencias, fuera del cribado PRISMA).} Fundamentan decisiones de
proceso o de diseño discutidas en otros capítulos, pero no son sistemas
de biblioteca ni estudios de rendimiento comparables entre sí:
ingeniería de requisitos (\citet{bano2019elicitation,
palomares2021elicitation, yu1997istar, kurtanovic2017classifying}),
prácticas ágiles (\citet{truong2025personality, rahman2024agile}),
diseño responsivo (\citet{parlakkilic2022responsive,
aienobe2025responsive}), seguridad de transporte
(\citet{mckay2019tls}), fundamentos de arquitectura de software
(\citet{sommerville2011se, fowler2002poeaa, walls2022springinaction,
bass2021softwarearch, brown2023c4model}) y muestreo en investigación
empírica de ingeniería de software (\citet{baltes2022sampling}, ya
citado en el \autoref{cap:amenazas-validez}). No se les aplicó un
embudo de selección porque no fueron evaluadas como candidatas
alternativas entre sí -- se incorporaron directamente por relevancia
temática con un capítulo o decisión concreta del proyecto.

\textbf{Nota de alcance -- referencias de JWT/autenticación reservadas
para Marco Teórico.} Tres referencias adicionales sobre el mecanismo
JWT (\citet{bucko2023jwtbehavior, shatnawi2024pythonjwt,
ahmed2019jwtscheme}, verificadas contra Crossref con el mismo rigor que
el resto del acervo) están en \texttt{docs/bibliografia.bib} pero
\textbf{no} se desarrollan en este capítulo: no son sistemas
bibliotecarios comparables (Grupo~A) ni fundamento metodológico general
(Grupo~B), sino evidencia técnica específica sobre JWT -- el mecanismo
de autenticación que SGB-SaaS ya implementa
(\autoref{cap:diseno-arquitectura}). Corresponden al capítulo de Marco
Teórico (\texttt{docs/capitulos/04-marco-teorico.tex}), que todavía no
se ha redactado; se deja esta nota para que, cuando se escriba, incluya
esas tres referencias en su discusión de JWT en vez de repetir aquí un
análisis que este capítulo no está en condiciones de sostener todavía.

\begin{figure}[h]
\centering
\begin{tikzpicture}[
    node distance=0.55cm and 1.6cm,
    box/.style={rectangle, draw=verdeuteq, thick, rounded corners,
      fill=grisclaro, text width=6.6cm, align=center, minimum height=1.1cm,
      font=\small},
    excl/.style={rectangle, draw=gapcolor, thick, rounded corners,
      fill=white, text width=5.2cm, align=left, minimum height=1.1cm,
      font=\footnotesize},
    arr/.style={-{Latex[length=2.2mm]}, thick, draw=verdeuteq}
  ]
  \node[box] (identificados)
    {\textbf{Identificados}\\[2pt] Grupo~A: 15 referencias\\
     (13 de curación previa del equipo + 2 de búsqueda dirigida
     Crossref/IEEE~Xplore para esta tarea)};
  \node[box, below=of identificados] (duplicados)
    {\textbf{Tras eliminar duplicados}\\[2pt] 15 referencias\\
     (0 duplicados encontrados)};
  \node[box, below=of duplicados] (cribados)
    {\textbf{Cribados por título/resumen}\\[2pt] 15 evaluados\\
     0 excluidos en esta fase};
  \node[excl, right=of cribados] (nota1)
    {Los 15 ya habían sido pre-filtrados por relevancia temática con
     SGB-SaaS antes de incorporarse al acervo -- por eso el cribado de
     título/resumen no excluye nada adicional aquí.};
  \node[box, below=of cribados] (textocompleto)
    {\textbf{Evaluados a texto completo / resumen extendido}\\[2pt]
     15 evaluados};
  \node[excl, right=of textocompleto] (nota2)
    {5 excluidos de la tabla comparativa primaria: son revisiones o
     artículos de posición que sintetizan el estado del arte
     (\citet{ayinde2026ai, liabor2023cataloging, adetayo2023chatgpt,
     yan2023chatbotsslr, ayemowa2024genai}), no sistemas primarios
     directamente comparables fila a fila. Se mantienen citados en la
     \autoref{sec:panorama-relacionados} como contexto.};
  \node[box, below=of textocompleto, fill=verdelight!25] (incluidos)
    {\textbf{Incluidos en la tabla comparativa de sistemas primarios}\\[2pt]
     \textbf{10 referencias} (\autoref{tab:trabajos-comparativa})};

  \draw[arr] (identificados) -- (duplicados);
  \draw[arr] (duplicados) -- (cribados);
  \draw[arr] (cribados) -- (textocompleto);
  \draw[arr] (textocompleto) -- (incluidos);
\end{tikzpicture}
\caption{Flujo de selección adaptado de PRISMA~2020 para el Grupo~A
  (dominio bibliotecario y arquitectura de acceso a datos comparable).
  El Grupo~B (15 referencias de apoyo metodológico, \autoref{sec:proceso-seleccion})
  no se somete a este embudo.}
\label{fig:prisma-flow}
\end{figure}

\section{Panorama de trabajos relacionados}
\label{sec:panorama-relacionados}

\subsection{Sistemas de gestión bibliotecaria: automatización, IA y
arquitectura de acceso a datos}

La literatura reciente sobre automatización bibliotecaria confirma una
tendencia clara hacia la incorporación de IA y de tecnologías de
identificación automática (QR, RFID), documentada en revisiones amplias
como \citet{ayinde2026ai} (adopción de IA en bibliotecas académicas) y
\citet{yan2023chatbotsslr} / \citet{ayemowa2024genai} (revisiones
sistemáticas específicas de chatbots y de sistemas de recomendación con
IA generativa, respectivamente). \citet{adetayo2023chatgpt} documenta,
a nivel editorial, la llegada de ChatGPT como punto de inflexión para
los chatbots de biblioteca, y \citet{liabor2023cataloging} describe
técnicas de catalogación digital sin evaluación empírica cuantitativa
propia. Ninguna de estas cinco es un sistema primario evaluable
fila a fila contra SGB-SaaS (\autoref{fig:prisma-flow}), por lo que se
citan aquí como contexto y no aparecen en la tabla comparativa.

La \autoref{tab:trabajos-comparativa} compara SGB-SaaS contra los 10
trabajos primarios seleccionados. Una nota de honestidad metodológica
aplica a varias filas: cuando el resumen indexado consultado no permitía
confirmar con certeza la pila tecnológica, el método de evaluación
empírica o las limitaciones declaradas de un trabajo -- por no haberse
podido acceder al texto completo del artículo, frecuentemente detrás de
un muro de pago editorial -- la celda correspondiente dice
explícitamente ``no confirmado a partir del resumen indexado
disponible'' en vez de completarse con una suposición razonable pero no
verificada.

\begin{landscape}
\begin{center}
\footnotesize
\renewcommand{\arraystretch}{1.15}
\setlength{\tabcolsep}{3pt}
\begin{longtable}{|p{2.1cm}|p{0.6cm}|p{2.6cm}|p{2.6cm}|p{2.9cm}|p{3.1cm}|p{2.9cm}|p{3.2cm}|}
\caption{Comparación de trabajos relacionados frente a SGB-SaaS}
\label{tab:trabajos-comparativa}\\
\hline
\rowcolor{verdeuteq}
\textcolor{white}{\textbf{Referencia}} & \textcolor{white}{\textbf{Año}} &
\textcolor{white}{\textbf{Dominio}} & \textcolor{white}{\textbf{Pila
tecnológica}} & \textcolor{white}{\textbf{Patrones arquitectónicos}} &
\textcolor{white}{\textbf{Evaluación empírica reportada}} &
\textcolor{white}{\textbf{Limitaciones declaradas}} &
\textcolor{white}{\textbf{Diferencia frente a SGB-SaaS}} \\
\hline
\endfirsthead
\hline
\rowcolor{verdeuteq}
\textcolor{white}{\textbf{Referencia}} & \textcolor{white}{\textbf{Año}} &
\textcolor{white}{\textbf{Dominio}} & \textcolor{white}{\textbf{Pila
tecnológica}} & \textcolor{white}{\textbf{Patrones arquitectónicos}} &
\textcolor{white}{\textbf{Evaluación empírica reportada}} &
\textcolor{white}{\textbf{Limitaciones declaradas}} &
\textcolor{white}{\textbf{Diferencia frente a SGB-SaaS}} \\
\hline
\endhead
\hline
\endfoot

\citet{supriyono2018qr} & 2018 &
Préstamo/devolución con QR en biblioteca escolar privada (Surakarta,
Indonesia) &
Web (Bootstrap) + MySQL + lector de cámara USB &
Monolito web tradicional, CRUD directo sobre MySQL (no se reporta ORM
ni SP) &
Pruebas funcionales/de aceptación con usuarios reales de la biblioteca;
sin métricas de carga ni estadística formal &
Caso único (una escuela); sin datos de escalabilidad &
SGB-SaaS usa QR solo como credencial de acceso (REQ-F-019), no como eje
central del préstamo, y aporta evidencia de carga formal (k6) ausente
aquí \\
\hline

\citet{mukund2021rfid} & 2021 &
LMS inteligente basado en RFID con analítica predictiva de demanda de
libros &
RFID + interfaz de escritorio PyQt &
Sistema de escritorio con lectura de hardware RFID; no arquitectura web
multi-tenant &
No confirmado a partir del resumen indexado disponible &
No confirmado a partir del resumen indexado disponible; dependencia
evidente de hardware RFID especializado &
SGB-SaaS es 100\,\% web multiusuario concurrente sin hardware
propietario, con auditoría OWASP y pruebas de carga documentadas \\
\hline

\citet{ng2024libbot} & 2024 &
Chatbot asistente bibliotecario (Lib-Bot) &
No confirmado a partir del resumen indexado disponible &
Asistente conversacional dedicado a consultas bibliotecarias &
No confirmado a partir del resumen indexado disponible &
No confirmado a partir del resumen indexado disponible &
SGB-SaaS integra un chatbot con IA generativa real (Gemini,
REQ-F-028) para recomendaciones \emph{y} consultas dentro del mismo
sistema transaccional, con \emph{rate limiting} propio en Redis \\
\hline

\citet{rajacic2025gdpr} & 2025 &
Validación de cumplimiento GDPR en un LMS real (BISIS, +60 bibliotecas
en Serbia) integrando el framework TeMDA &
BISIS (LMS en producción) + TeMDA (framework de validación dirigido por
modelos) &
Integración de un framework externo de validación de privacidad sobre
un LMS ya productivo; caso de estudio con diario de desarrollo &
Estudio de caso cualitativo (diario de desarrollo); sin métricas
cuantitativas de rendimiento ni pruebas de carga &
Caso único, cualitativo, sin generalización estadística &
SGB-SaaS \textbf{no} implementa un framework dedicado de validación
GDPR -- gap real, no solo diferencia --; solo cuenta con bitácora de
auditoría interna (REQ-F-003) sin verificación normativa formal \\
\hline

\citet{wayesa2023recsys} & 2023 &
Recomendación híbrida de libros basada en patrones y relaciones
semánticas &
No confirmado en detalle a partir del resumen indexado disponible &
Módulo de recomendación independiente, no integrado a un LMS
transaccional completo &
Evaluación de calidad de recomendación (según título/resumen); detalle
cuantitativo no confirmado en el resumen consultado &
No confirmado en detalle a partir del resumen indexado disponible &
SGB-SaaS no implementa un modelo de recomendación semántica dedicado;
delega la recomendación a un chatbot con IA generativa de propósito
general (Gemini) -- \emph{trade-off} explícito, no superioridad \\
\hline

\citet{hu2025bigdata} & 2025 &
Analítica de big data e IA en servicios de biblioteca: recomendación
personalizada y optimización de catálogo &
Framework de IA con autoencoder variacional (VAE) y optimizador Adam
(según fragmento de resumen disponible) &
Pipeline de analítica/ML sobre datos de catálogo, orientado a
optimización, no a transacciones de préstamo &
No confirmado en detalle -- acceso completo restringido por muro de
pago editorial (IEEE Access) &
No confirmado en detalle (mismo motivo) &
SGB-SaaS no entrena un modelo de ML propio (VAE ni similar); delega la
personalización a una API de IA generativa externa -- arquitectura más
simple, dependiente de proveedor externo \\
\hline

\citet{ehrenpreis2022chatbot} & 2022 &
Implementación de un chatbot propietario (Ivy) en el sitio web de una
biblioteca universitaria (Lehman College) -- primer caso documentado en
biblioteca académica &
Software de chatbot educativo propietario embebido en sitio web
institucional &
Chatbot de terceros integrado solo al sitio informativo, sin conexión a
las transacciones del sistema de préstamos &
Estudio de caso descriptivo de implementación; sin métricas
cuantitativas de precisión ni de carga &
Descripción de un solo caso; sin datos cuantitativos de efectividad del
chatbot &
El chatbot de SGB-SaaS no es un producto de terceros aislado: es un
módulo (REQ-F-028) integrado a los datos transaccionales reales del
catálogo \\
\hline

\citet{shivaji2024intelligentlms} & 2024 &
Sistema inteligente de gestión bibliotecaria con RFID, módulo GSM y
microcontrolador AVR &
RFID + GSM + microcontrolador AVR + Visual Basic (según palabras clave
del resumen indexado) &
Sistema embebido/de escritorio orientado a automatización física
(notificación por SMS vía GSM), no arquitectura SaaS web &
No confirmado en detalle a partir del resumen indexado disponible &
No confirmado en detalle (mismo motivo); dependencia de hardware
especializado evidente por el propio diseño &
SGB-SaaS es una plataforma SaaS 100\,\% software, desplegable en
cualquier navegador, con notificación por correo (SMTP) en vez de
SMS/GSM \\
\hline

\citet{colley2018ormperformance} & 2018 &
Impacto de frameworks ORM en el rendimiento de consultas relacionales
(estudio general de acceso a datos, no específico de bibliotecas) &
Stack de aplicación Java líder de mercado (no se especifica dominio
bibliotecario) &
Identifica patrones de rendimiento indeseables (anti-patrones)
producidos por el uso de ORM frente a SQL optimizado &
Comparación empírica real de tiempos de ejecución y uso de memoria entre
configuraciones ORM &
No evalúa explícitamente una arquitectura híbrida ORM+procedimientos
almacenados como estrategia de mitigación &
SGB-SaaS \textbf{sí} implementa y mide empíricamente (k6, Wilcoxon) la
estrategia híbrida CRUD-ORM/SP que este trabajo identifica como
necesaria pero no llega a ejecutar ni medir (\autoref{sec:sp-invocacion}) \\
\hline

\citet{kusuma2017rediscaching} & 2017 &
Comparación de estrategias de caché (acelerador HTTP + Redis) en
WordPress multisitio -- no bibliotecario &
WordPress + Varnish (caché HTTP) + Redis (caché de objetos) + MySQL &
\emph{Cache-aside} separado por tipo de objeto: estático vía Varnish,
consultas de BD vía Redis &
Medición real de uso de CPU de MySQL ($-25$\,\%) y tráfico de red
($-94$\,\%) bajo carga, con caché Redis activo &
Contexto específico de WordPress multisitio; sin test estadístico
formal ($p$-valor, tamaño de efecto) &
SGB-SaaS usa un patrón \emph{cache-aside} equivalente
(\autoref{sec:cache-redis}) sobre una API REST propia, con evaluación
estadística formal (Wilcoxon, Cliff's delta) que este trabajo no reporta \\
\hline
\end{longtable}
\end{center}
\end{landscape}

\subsection{Ingeniería de requisitos: elicitación y clasificación}

La calidad del proceso de elicitación de requisitos es un tema activo de
investigación empírica: \citet{bano2019elicitation} estudian, mediante un
análisis de errores comunes, cómo se enseña y se aprende a conducir
entrevistas de elicitación, y \citet{palomares2021elicitation} amplían el
panorama con un estudio de estado de la práctica en 12 empresas reales.
\citet{yu1997istar} aporta el marco conceptual clásico ($i^{*}$) para
razonar sobre requisitos en fase temprana en términos de actores e
intenciones, y \citet{kurtanovic2017classifying} atacan un problema
técnico directamente relevante para este proyecto: la clasificación
automática de requisitos funcionales frente a no funcionales mediante
aprendizaje supervisado -- la misma distinción F/NF que estructura la
matriz de trazabilidad de SGB-SaaS (\autoref{cap:ingenieria-requisitos}),
aunque en este proyecto esa clasificación se hizo manualmente y no con un
clasificador entrenado.

\subsection{Prácticas ágiles y efectividad de equipo}

\citet{truong2025personality} encuentran, en equipos Scrum de empresas
vietnamitas, una asociación medible entre rasgos de personalidad del
equipo y su efectividad percibida, y \citet{rahman2024agile} reportan
evidencia de que la adopción de prácticas ágiles impacta la
productividad de equipos de desarrollo. Ambos trabajos son relevantes
como marco de referencia para contextualizar (no medir formalmente,
dado que está fuera del alcance de este proyecto) el propio proceso de
un equipo de 3 personas sin dedicación exclusiva que construyó SGB-SaaS
(\autoref{cap:amenazas-validez}).

\subsection{Diseño responsivo de interfaces web}

\citet{parlakkilic2022responsive} evalúan el efecto del diseño
responsivo sobre la usabilidad de sitios web académicos durante la
pandemia, y \citet{aienobe2025responsive} revisan patrones de diseño
responsivo orientados a mejorar UX/UI. Ambos fundamentan la decisión de
construir el frontend de SGB-SaaS como una aplicación Angular
responsiva -- decisión cuya validación de usabilidad real (SUS) sigue
pendiente al momento de escribir este documento
(\autoref{sec:validez-externa}).

\subsection{Seguridad de transporte}

\citet{mckay2019tls} (NIST SP~800-52 Rev.~2) documentan las guías
oficiales para la selección y configuración de implementaciones TLS, la
referencia normativa detrás de las decisiones de transporte seguro de
credenciales de SGB-SaaS (cookies \texttt{HttpOnly}, \autoref{cap:implementacion}).

\subsection{Fundamentos de arquitectura y desarrollo de software}

Las decisiones de diseño documentadas en el \autoref{cap:diseno-arquitectura}
y el \autoref{cap:implementacion} se apoyan en literatura fundacional
consolidada: \citet{sommerville2011se} para el proceso de ingeniería de
software en general, \citet{fowler2002poeaa} para los patrones de
arquitectura de aplicaciones empresariales (incluyendo el patrón de
acceso a datos que motiva la estrategia híbrida ORM/SP),
\citet{walls2022springinaction} para el framework Spring usado en el
backend, \citet{bass2021softwarearch} para el vocabulario de atributos
de calidad usado en el mapeo a ISO/IEC~25010
(\autoref{cap:diseno-arquitectura}), y \citet{brown2023c4model} para la
notación C4 usada para documentar la arquitectura -- citado, como se
aclaró en la \autoref{sec:estrategia-busqueda}, como referencia técnica
normativa y no como evidencia empírica de investigación.

\section{Brecha identificada}
\label{sec:brecha-identificada}

Ningún trabajo revisado en este capítulo integra, en un solo sistema,
las cuatro características que SGB-SaaS combina: (i)~autenticación
basada en JWT con control de acceso por roles (RBAC); (ii)~un catálogo
con autocompletado de ISBN; (iii)~un chatbot con IA generativa real --no
un motor de intents fijos-- usado simultáneamente para recomendaciones y
para consultas sobre el catálogo (REQ-F-028); y (iv)~una arquitectura
híbrida de acceso a datos (CRUD elemental vía ORM, operaciones
multi-tabla vía procedimientos almacenados) respaldada por evidencia
empírica de rendimiento reproducible (k6, 5 corridas independientes,
prueba de Wilcoxon pareada, tamaño de efecto de Cliff's delta).

De los 10 trabajos primarios comparados, los más cercanos a SGB-SaaS en
cada dimensión individual son: \citet{ng2024libbot} y
\citet{ehrenpreis2022chatbot} en la dimensión de chatbot (pero ninguno
de los dos reporta el uso de un modelo de IA generativa de propósito
general integrado a las transacciones del sistema, a diferencia de
Gemini en SGB-SaaS); \citet{colley2018ormperformance} en la dimensión de
arquitectura de acceso a datos (identifica el problema que la estrategia
híbrida de SGB-SaaS resuelve, pero no implementa ni mide esa estrategia);
y \citet{kusuma2017rediscaching} en la dimensión de caché (mide el
efecto de Redis, pero sin la evaluación estadística formal que
\texttt{docs/mediciones/perf/REPORT.md} aporta para SGB-SaaS). Ningún
trabajo del acervo revisado combina las cuatro dimensiones en un solo
sistema evaluado empíricamente.

Esta afirmación de brecha se hace con dos advertencias explícitas de
honestidad metodológica, coherentes con el resto de este documento:
primera, está acotada estrictamente a las 30 referencias efectivamente
revisadas en este capítulo (\autoref{sec:estrategia-busqueda}) -- no es
una afirmación sobre la totalidad de la literatura publicada sobre
sistemas de gestión bibliotecaria, que es considerablemente más extensa
que el acervo aquí reunido; segunda, la brecha se refiere a la
\emph{combinación} de las cuatro características, no a la novedad de
cada característica individual, varias de las cuales (chatbots,
recomendación, RFID, QR) están bien establecidas por separado en la
literatura revisada. Esta misma limitación de alcance de la revisión se
retoma como amenaza a la validez externa en el
\autoref{cap:amenazas-validez}.
